% Vietnamese translation for OI Public Library
% Source-File: IOI中国国家候选队论文/2003/林希德/林希德.ppt
\documentclass[11pt]{article}
\usepackage{oipl-vn}

\title{Bản trình chiếu: Tìm chuỗi con lặp lớn nhất}
\author{Lin Xide}
\date{2003}

\begin{document}
\maketitle

\origtitle{Tìm chuỗi con lặp lớn nhất}
\sourcepath{IOI China candidate team papers / 2003 / Lin Xide / Lin Xide.ppt}

\section{Mục tiêu của bài trình bày}

Bộ slide trình bày bài toán tìm chuỗi con lặp lớn nhất trong một chuỗi chữ
cái in hoa. Nhiều trang gốc chứa hình ảnh hoặc đối tượng nhúng, nên bản ghi
chú này tập trung vào nội dung thuật toán đã được bài viết đi kèm giải thích:
định nghĩa chu kỳ, cây hậu tố, mở rộng KMP và cách kết hợp hai công cụ này để
đạt thuật toán tuyến tính.

Các đoạn mã dài hoặc đối tượng nhúng trong nguồn gốc không được chép lại
nguyên văn. Theo ngoại lệ về mẫu mã, bản dịch chỉ tóm tắt vai trò của chúng:
cây hậu tố dùng để phân rã chuỗi, còn mở rộng KMP dùng để tính nhanh các độ
dài khớp cần thiết.

\section{Định nghĩa bài toán}

Với chuỗi \(S\), một số nguyên dương \(p\) là chu kỳ của \(S\) nếu
\[
  1\le x\le |S|-p \quad\Longrightarrow\quad S_x=S_{x+p}.
\]
Khi đó
\[
  e=\frac{|S|}{p}
\]
là số mũ của \(S\). Nếu \(e\ge2\), ta gọi \(S\) là một chuỗi lặp kích thước
\(p\). Bài toán đặt ra: với chuỗi \(W\) độ dài
\[
  1\le n\le10^5,
\]
tìm trong mọi chuỗi con của \(W\) một chuỗi lặp có chu kỳ lớn nhất.

Ký hiệu
\[
  W(u,v)
\]
là chuỗi con bắt đầu tại vị trí \(u\) và kết thúc tại vị trí \(v\). Cách trực
tiếp là duyệt chu kỳ \(p\) từ lớn xuống nhỏ rồi kiểm tra có hai khối liên tiếp
độ dài \(p\) giống nhau hay không. Cách này dễ hiểu nhưng tốn \(O(n^2)\).

\section{Cây hậu tố trong slide}

Để xử lý mọi hậu tố trong thời gian tuyến tính, ta thêm ký tự kết thúc đặc
biệt:
\[
  W\leftarrow W\$,
\]
với \(\$\) không xuất hiện trong chuỗi ban đầu. Cây hậu tố là trie nén của
tất cả hậu tố \(W(i,n)\). Mỗi cạnh \(E(u,v)\) được ghi bằng cặp chỉ số
\[
  E(u,v)=(\ell,r),
\]
biểu diễn nhãn cạnh là chuỗi con \(W(\ell,r)\).

Nếu nối các nhãn trên đường từ gốc tới nút \(x\), ta được chuỗi \(S_x\). Mỗi
hậu tố \(W(i,n)\) ứng với một lá \(\mathit{Leaf}_i\). Khi chèn hậu tố thứ
\(i\), đặt \(\mathit{Head}_i\) là tiền tố chung dài nhất giữa \(W(i,n)\) và
một hậu tố đã chèn trước đó; phần còn lại là \(\mathit{Tail}_i\). Cây hậu tố
được cập nhật bằng cách tìm vị trí của \(\mathit{Head}_i\), tách cạnh nếu cần,
rồi thêm lá cho \(\mathit{Tail}_i\).

Slide nhắc ý tưởng liên kết hậu tố của McCreight. Nếu
\[
  \mathit{Head}_{i-1}=az,
\]
thì khi chuyển sang hậu tố kế tiếp, chuỗi \(z\) đã biết là một tiền tố xuất
hiện trước đó. Liên kết hậu tố từ nút \(az\) sang nút \(z\) cho phép tìm
\(\mathit{Head}_i\) mà không quay lại gốc. Đại lượng
\[
  i+|\mathit{Head}_i|
\]
không giảm, nên tổng số ký tự được duyệt là \(O(n)\).

\section{KMP và mở rộng KMP}

KMP cổ điển dùng hàm tiền tố để tránh so sánh lại những ký tự đã biết. Trong
bài toán này, ta cần nhiều giá trị độ dài khớp chứ không chỉ cần biết mẫu có
xuất hiện hay không. Vì vậy slide dùng dạng mở rộng KMP, tương đương với
việc tính các giá trị LCP kiểu Z-function.

Với mẫu \(T\) và chuỗi \(S\), định nghĩa
\[
  B_i=\operatorname{lcp}\bigl(T,S(i,n)\bigr).
\]
Trước hết tính
\[
  A_i=\operatorname{lcp}\bigl(T,T(i,m)\bigr),
\]
rồi dùng đoạn khớp xa nhất đã biết để suy ra hoặc mở rộng từng \(B_i\). Vì
biên phải của đoạn khớp chỉ tăng, tổng số lần so sánh ký tự là tuyến tính.

Trong thuật toán chính, mở rộng KMP cung cấp nhanh hai loại thông tin:
khớp tiền tố khi mở một chuỗi sinh sang phải, và khớp hậu tố khi mở sang trái
trên chuỗi đảo.

\section{Chuỗi con tối ưu}

Với chuỗi con \(S=W(u,v)\) và chu kỳ \(L\), ta gọi \(S\) là chuỗi con tối ưu
có chu kỳ \(L\) nếu:
\begin{itemize}
  \item \(S\) có chu kỳ \(L\);
  \item không thể mở rộng \(S\) sang trái mà vẫn giữ chu kỳ \(L\);
  \item không thể mở rộng \(S\) sang phải mà vẫn giữ chu kỳ \(L\).
\end{itemize}
Nếu liệt kê được các chuỗi con tối ưu và chu kỳ của chúng, đáp án là chuỗi
có chu kỳ lớn nhất. Cách kiểm tra một chu kỳ \(L\) là chọn một chuỗi sinh độ
dài \(L\), mở rộng sang hai phía đến khi không còn khớp, rồi xem độ dài sau
mở rộng có đạt ít nhất \(2L\) hay không.

\section{Phân rã chuỗi}

Slide phân rã
\[
  W=U_1+U_2+\cdots+U_m
\]
theo quy tắc: nếu phần đã xử lý là
\[
  P=U_1+\cdots+U_{i-1},\qquad W=P+Q,
\]
thì nếu ký tự đầu \(Q_1\) chưa xuất hiện trong \(P\), đặt \(U_i=Q_1\); ngược
lại, \(U_i\) là tiền tố dài nhất của \(Q\) đã xuất hiện trong \(P\).

Ví dụ trong bài viết:
\[
  W=\texttt{ABAABABAABAAB}
\]
có phân rã
\[
  \texttt{|A|B|A|AB|ABAAB|AAB|}.
\]
Phân rã này được tính bằng cây hậu tố, vì mỗi \(U_i\) chính là một ký tự mới
hoặc \(\mathit{Head}_{|P|+1}\).

\section{Giới hạn nơi có thể có nghiệm}

Vai trò của phân rã là giới hạn vị trí của chuỗi con tối ưu. Xét một chuỗi
con tối ưu \(S\) có điểm kết thúc nằm trong mảnh \(U_i\). Nếu \(S\) bắt đầu
trong chính \(U_i\), nó đã xuất hiện trong phần trước theo định nghĩa của
\(U_i\), nên sẽ được xét ở bước trước đó. Trường hợp cần xét là \(S\) bắt đầu
trước \(U_i\) và kết thúc trong \(U_i\).

Đặt \(U=U_i\), và đặt \(V\) là hậu tố của phần trước đó
\[
  P=U_1+\cdots+U_{i-1}
\]
có độ dài
\[
  |U|+2|U_{i-1}|.
\]
Khi ấy mọi chuỗi con tối ưu kết thúc trong \(U\) đều có điểm bắt đầu nằm
trong \(V\). Nhờ vậy, thay vì xét toàn bộ chuỗi, ta chỉ cần xét các bài toán
cục bộ trên \(V+U\).

\section{Tìm chuỗi sinh trong từng cục bộ}

Với một cặp \(V,U\), xét chu kỳ \(L\). Nếu chuỗi con tối ưu bắt đầu trong
\(V\), kết thúc trong \(U\), và có độ dài ít nhất \(2L\), thì ít nhất một
trong hai phần của nó nằm trong \(V\) hoặc \(U\) có độ dài không nhỏ hơn
\(L\). Hai trường hợp đối xứng; xét trường hợp phần nằm trong \(U\) đủ dài.
Khi đó
\[
  U(1,L)
\]
có thể được dùng làm chuỗi sinh, rồi mở rộng sang trái và sang phải.

Các đại lượng phụ trợ được slide dùng là
\[
  \mathit{Lp}_L=\operatorname{lcp}\bigl(U,U(L+1,|U|)\bigr)
\]
và
\[
  \mathit{Ls}_L=\operatorname{lcs}\bigl(V,V+U(1,L)\bigr).
\]
Ở đây \(\operatorname{lcs}\) là độ dài hậu tố chung dài nhất. Toàn bộ các
giá trị \(\mathit{Lp}_L\) và \(\mathit{Ls}_L\) được tính bằng mở rộng KMP, nên
chi phí bình quân cho mỗi \(L\) là \(O(1)\).

\section{Thuật toán tổng quát}

Quy trình của bài trình bày có thể tóm tắt như sau:
\begin{enumerate}
  \item Dựng cây hậu tố và phân rã \(W\) thành \(U_1,\ldots,U_m\).
  \item Với mỗi \(i\ge2\), đặt \(U=U_i\), chọn \(V\) là hậu tố của phần đã xử
        lý có độ dài \(|U|+2|U_{i-1}|\).
  \item Tính các mảng mở rộng cho cặp \(V,U\).
  \item Với mọi \(1\le L\le |U|\), dùng \(U(1,L)\) làm chuỗi sinh, mở rộng
        bằng \(\mathit{Lp}_L\) và \(\mathit{Ls}_L\), rồi kiểm tra độ dài sau
        mở rộng có ít nhất \(2L\) hay không.
  \item Làm tương tự cho trường hợp đối xứng khi phần dài nằm trong \(V\).
  \item Cập nhật đáp án bằng chu kỳ lớn nhất tìm được.
\end{enumerate}

Tổng độ dài các mảnh \(U_i\) là \(O(n)\), còn cách chọn \(V\) bảo đảm tổng
chi phí của mọi bài toán cục bộ cũng tuyến tính. Vì vậy toàn bộ thuật toán
chạy trong \(O(n)\) thời gian và dùng \(O(n)\) bộ nhớ.

\section{Tổng kết}

Bài trình bày không chỉ giới thiệu cây hậu tố hoặc KMP riêng lẻ. Điểm chính
là cách phối hợp chúng: cây hậu tố tạo ra một phân rã làm hẹp vùng có thể chứa
nghiệm, còn mở rộng KMP cho phép kiểm tra tất cả chu kỳ trong từng vùng với
tổng chi phí tuyến tính. Nhờ đó một bài toán nhìn ban đầu giống tìm kiếm
mọi chuỗi con \(O(n^2)\) được đưa về thuật toán \(O(n)\).

\end{document}
